% DRAFT (not \input in main.tex until the framing is decided). Numbers are filled from results/t1/g*_*.json.
\subsection{The observable geometry of the behavioural state}
\label{sec:geometry}
\Cref{prop:gauge} says that the hidden coordinates are not observable. Their \emph{geometry as seen through the outputs} is. Let $P_x=\bigotimes_{j}\mathrm{Bernoulli}\big(\Psi_{q_j}(x)\big)$ be the law of the answers after context $x$, and let $d_Q(x,x')$ be the Fisher--Rao geodesic distance between $P_x$ and $P_{x'}$.
\begin{theorem}[observable geometry]\label{thm:geometry}
(a) $d_Q(x,x')^2=\sum_j\big(\vartheta_j(x)-\vartheta_j(x')\big)^2$ with $\vartheta_j=2\arcsin\sqrt{\Psi_{q_j}}$: the angles embed the answer laws isometrically in Euclidean space.
(b) $d_Q$ depends on the outputs only, so it is invariant under the gauge group of \cref{prop:gauge}; up to a constant factor, the Fisher metric is the only Riemannian metric on the answer laws that is invariant under sufficient statistics (\v{C}encov).
(c) Near a state $s$, $d_Q^2=\Delta^{\top}G_Q(s)\Delta+O(|\Delta|^3)$ with $G_Q=\sum_j\Psi_j(1-\Psi_j)\nabla a_j\nabla a_j^{\top}$ for $\Psi_j=\sigma(a_j(s))$: the observable distance is the hidden distance in the pulled-back Fisher metric, which gains rank as questions are added.
(d) With $n$ sampled answers per question, let $\mu_n(p)=\E\,\hat\vartheta$ for the Anscombe angle $\hat\vartheta=2\arcsin\sqrt{(k+3/8)/(n+3/4)}$, $k\sim\mathrm{Bin}(n,p)$: a known, strictly increasing function with $\mu_n\to\vartheta$. For two independent halves $A,B$ of the samples, $\hat D^2(x,x')=\sum_j(\hat\vartheta^A_j(x)-\hat\vartheta^A_j(x'))(\hat\vartheta^B_j(x)-\hat\vartheta^B_j(x'))$ has $\E\hat D^2=\sum_j\big(\mu_n(\Psi_j(x))-\mu_n(\Psi_j(x'))\big)^2$ exactly, for every $\Psi$; classical scaling then returns the coordinates up to a rigid motion.
\end{theorem}
No model is trained: (a) is a change of variables, and (d) costs $m\cdot n$ calls per context on any chat API. Part (c) gives two predictions that we test: agreement with the hidden geometry rises with $m$ and saturates, and it peaks where the answers are read out (middle to late layers).

Where full log-probabilities are available, the final state is not estimated but recovered.
\begin{theorem}[exact recovery]\label{thm:exact}
Let the logits be $Wu$ with $u=x/\mathrm{rms}(x)$, the normalisation gain folded into $W$, and no output bias. The centred log-probabilities $c(x)$ lie in a $d$-dimensional subspace; write $c=Q_d\,b$. Then $u=Kb$ for an unknown invertible $K$, and $\|u\|^2=d$ gives $b_i^{\top}Sb_i=d$ with $S=K^{\top}K$, a linear system in $S$. If the $b_ib_i^{\top}$ span the symmetric matrices (generically, $N\ge d(d+1)/2$ contexts), $S$ is unique and the Gram matrix $U^{\top}U=B^{\top}SB$ of the final states is exact; the states follow up to an orthogonal map, which is the gauge group.
\end{theorem}
Below $d(d+1)/2$ contexts, $S$ is not determined; above, it is. The theory thus predicts a sharp transition at a known count.

% ---------------- appendix proofs (DRAFT)
\begin{proof}[Proof of \cref{thm:geometry}]
(a) The Fisher information of $\mathrm{Bernoulli}(p)$ is $1/(p(1-p))$ and $\mathrm{d}\vartheta/\mathrm{d}p=1/\sqrt{p(1-p)}$, so the metric is $\mathrm{d}\vartheta^2$. The Fisher metric of a product of independent factors is the sum of the factors' metrics, so in $\vartheta$ the family is the box $[0,\pi]^m$ with the Euclidean metric. The box is convex, so geodesics are segments and the distance is Euclidean. Check for one factor: the Bhattacharyya coefficient of $\mathrm{Bernoulli}(\sin^2\alpha)$ and $\mathrm{Bernoulli}(\sin^2\beta)$ is $\cos(\alpha-\beta)$, and $2\arccos\cos(\alpha-\beta)=|\vartheta-\vartheta'|$.
(b) Invariance: $d_Q$ is a function of $(P_x,P_{x'})$ only, and \cref{prop:gauge} leaves every output law unchanged. Uniqueness is \v{C}encov's theorem for finite sample spaces (Campbell's extension to the positive cone; each answer law lives on $\{0,1\}^m$).
(c) $\mathrm{d}\vartheta_j=\Psi_j'/\sqrt{\Psi_j(1-\Psi_j)}$ and $\Psi_j'=\Psi_j(1-\Psi_j)\nabla a_j^{\top}\Delta$, so $\mathrm{d}\vartheta_j=\sqrt{\Psi_j(1-\Psi_j)}\,\nabla a_j^{\top}\Delta$; summing the squares gives $\Delta^{\top}G_Q\Delta$, and the remainder is third order by Taylor's theorem.
(d) $\hat\vartheta$ is increasing in $k$ and $k$ is stochastically increasing in $p$, so $\mu_n$ is strictly increasing; $\mu_n\to\vartheta$ by the law of large numbers and continuity. The halves are independent, so $\E[(a_A-a'_A)(a_B-a'_B)]=(\E a_A-\E a'_A)(\E a_B-\E a'_B)=(\mu_n(\Psi)-\mu_n(\Psi'))^2$ for each question; sum over questions. Classical scaling recovers a Euclidean configuration from its distance matrix up to a rigid motion (Young--Householder; Torgerson). Note that the familiar large-$n$ variance $1/(n+\frac12)$ of the Anscombe angle fails near $p\in\{0,1\}$ and its bias is $\cot\vartheta/(4n)$, so subtracting a variance term is not exact; the cross-fit needs neither.
\end{proof}
\begin{proof}[Proof of \cref{thm:exact}]
With $P=I-\mathbf{1}\mathbf{1}^{\top}/V$, the log-softmax removes a constant per context, so $c=PWu$ and every $c$ lies in the column space of $PW$, of dimension $d$ when $PW$ has full column rank. Let $Q_d$ be an orthonormal basis of it and $b=Q_d^{\top}c=Mu$ with $M=Q_d^{\top}PW$ invertible; then $u=Kb$ with $K=M^{-1}$. $\|u\|^2=b^{\top}K^{\top}Kb=d\,\overline{x^2}/(\overline{x^2}+\epsilon)$, which is $d$ up to the normalisation's $\epsilon$. The map $S\mapsto(b_i^{\top}Sb_i)_i$ is linear on the $d(d+1)/2$-dimensional space of symmetric matrices; it is injective iff the $b_ib_i^{\top}$ span that space, which for $b_i$ in general position needs and generically holds with $N\ge d(d+1)/2$. Then $S=K^{\top}K$ is the unique solution and $B^{\top}SB=U^{\top}U$. Any $L$ with $L^{\top}L=S$ satisfies $L=OK$ for an orthogonal $O$, so the states are fixed up to $O$, the group of \cref{prop:gauge}.
\end{proof}
